\section{Projected curvature and multiplicity projections}\label{sec:curvature}

All estimates in this section hold on the product of a fixed small
parameter neighborhood and a fixed small ball of coefficient arrays.
The constants are independent of any positive finite eigenvalue.
For each $s$ consider pointwise trace-free, $h_{0,s}$-Hermitian
endomorphisms of $G$. Let $P_s$ be their $L^2$ projection onto the
entire parallel algebra $\hh$, and put $Q_s=1-P_s$.
For $u\in Q_s$ define
\begin{equation}\label{eq:conjugatedcurvature}
\widehat K(s,Z,u)=e^{u/2}
\bigl(i\Lambda_sF_{\dbar_{s,Z},h_{0,s}e^u}\bigr)e^{-u/2}
-\kappa_E(s)I.
\end{equation}
This is Hermitian for $h_{0,s}$. Its pointwise trace vanishes:
the determinant holomorphic structure is the split determinant and
$\tr u=0$, so the determinant curvature is unchanged and has constant
contraction $R_E\kappa_E$.

\begin{lemma}\label{lem:complement}
The equation $Q_s\widehat K(s,Z,u)=0$ has a unique small solution
$u(s,Z)\in Q_s$, smooth in the finite parameters. For every fixed
nonnegative integer $a$,
\begin{equation}\label{eq:uquadratic}
\norm{u(s,Z)}_{C^a}\le C_aq_s(Z).
\end{equation}
It satisfies $u(s,k\cdot Z)=ku(s,Z)k^{-1}$ for $k\in K$.
\end{lemma}

\begin{proof}
At $Z=u=0$ the linearization of \eqref{eq:conjugatedcurvature} in
the real Hermitian metric variable is
\begin{equation}\label{eq:linearization}
L_s=\tfrac12(\Delta_{\dbar,s}+\Delta_{\partial,s}).
\end{equation}
The conjugation in \eqref{eq:conjugatedcurvature} is included in
this formula; it gives the symmetric real linearization even when
the split summands have different mean curvatures. The Kähler
identities give \eqref{eq:linearization}, and
\[
\langle u,L_su\rangle_{L^2}
=\tfrac12\bigl(\|\dbar u\|_{L^2}^2+\|\partial u\|_{L^2}^2\bigr).
\]
Its kernel is the parallel algebra. At $s=0$ this follows from
stability and vanishing cross-Hom groups. For nearby $s$, a zero
of the displayed sum is parallel for the same split connection;
there are no additional parallel maps. Conversely every
multiplicity block is parallel. The positive spectral gap of
$L_0$ on the complement therefore persists uniformly in a small
neighborhood, and so does its elliptic inverse.

Take an integer $\ell>d$, so that the relevant Sobolev spaces are
algebras on the real $2d$-dimensional manifold. The projected equation
is a smooth map from $Q_sH^{\ell+2}$ to $Q_sH^\ell$.
Identify these varying spaces smoothly with those at zero. The
parameter-dependent implicit-function theorem applies because the
derivative $Q_sL_sQ_s$ is uniformly invertible. At $Z=0$, the split
curvature is parallel; hence $u(s,0)=0$.

The first derivative of the contracted curvature in $Z$ vanishes.
Indeed its off-diagonal linear terms are contractions of covariant
derivatives of the harmonic forms $\beta_{ij,a,s}$; the Kähler
identities identify them with their vanishing Dolbeault adjoints.
Thus $D_Zu(s,0)=0$. Taylor expansion gives the quadratic bound in
Sobolev norm. Elliptic regularity gives \eqref{eq:uquadratic} for
each fixed spatial order. The group $K$ preserves the equation,
the entire parallel algebra, and its orthogonal complement.
Uniqueness proves the equivariance. It is unnecessary, and generally
false, that $K$ fix every vector of the parallel algebra.
\end{proof}

Write $\mathcal F_v(s,Z)$ for the volume-normalized integral of the
partial trace over $F_v$ of the $v,v$ block of
$\widehat K(s,Z,u(s,Z))$. Then
\begin{equation}\label{eq:reducedprojection}
\bigl(P_s\widehat K\bigr)_v
=I_{F_v}\otimes\frac{\mathcal F_v(s,Z)}{r_v}.
\end{equation}
In particular $\mathcal F\in\mathcal W$, and the full HE equation
is equivalent to $\mathcal F=0$ once the complement is solved.

\begin{proposition}\label{prop:projection}
There is a smooth residual $\mathcal R(s,Z)\in\mathcal W$ such that
\begin{equation}\label{eq:reducedexpansion}
\mathcal F(s,Z)=-b(s)+\mu_s(Z)+\mathcal R(s,Z),\qquad
\|D_Z^2\mathcal R(s,Z)\|\le Cq_s(Z).
\end{equation}
For every tuple $\Pi=(\Pi_v)$ of orthogonal multiplicity projections,
\begin{equation}\label{eq:projectionestimate}
\left|\sum_v\tr\bigl(\Pi_v\mathcal R_v(s,Z)\bigr)\right|
\le Cq_s(Z)q_s([\Pi,Z]).
\end{equation}
One constant works for all projections and all sufficiently small
$(s,Z)$, including projections of different ranks at different vertices.
\end{proposition}

The last factor measures the squared norm of precisely those arrow
blocks crossing between $\operatorname{im}\Pi$ and its orthogonal
complement. This factor, rather than a norm of $\Pi$ alone, will
retain control at an arbitrarily weak finite direction.

\begin{proof}
At $Z=0$, \eqref{eq:reducedprojection} gives $\mathcal F=-b$.
Harmonicity kills the linear term, as in Lemma~\ref{lem:complement}.
The direct quadratic term in the Chern curvature is the incoming-minus-
outgoing expression \eqref{eq:momentcomponents}, with its current
harmonic Gram form. The quadratic metric correction contributes
$L_su_2$, where $u_2$ is the quadratic Taylor term of $u$. Its
pairing with every parallel block vanishes, because $L_s$ is
self-adjoint and annihilates those blocks. This proves the constant,
linear, and quadratic parts of \eqref{eq:reducedexpansion}.

The unitary which equals $-I$ at all targets and $I$ at all sources
sends $Z$ to $-Z$ and acts trivially by conjugation on the parallel
algebra. Equivariance yields $\mathcal F(s,-Z)=\mathcal F(s,Z)$.
Consequently the cubic term also vanishes. Joint Taylor expansion,
with the background parameter restricted to a relatively compact
neighborhood, gives the Hessian bound in \eqref{eq:reducedexpansion}.
In particular it holds for each nearby $s$, with no comparison
between $|s|$ and $q_s(Z)$.

Fix $\Pi$ and decompose each arrow into its preserving and crossing
parts, writing $Z=A+C$. Explicitly,
\[
A_{ij,a}=\Pi_iZ_{ij,a}\Pi_j+(1-\Pi_i)Z_{ij,a}(1-\Pi_j),
\quad C=Z-A.
\]
The two parts are orthogonal in the harmonic Gram norm, and
$q_s(C)=q_s([\Pi,Z])$. This remains true for nondiagonal arrow
Gram matrices, since preserving and crossing matrix blocks are
orthogonal for each pair of arrow indices.

At $C=0$ the harmonic bundle is a holomorphic direct sum
$G_\Pi\oplus G_{1-\Pi}$. A circle acting scalarly on $G_\Pi$
and trivially on its complement fixes $A$. Equivariance and
uniqueness of $u$ force $u(s,A)$ to preserve this sum. Integrated
Chern--Weil trace on the first summand is its degree, whence
\begin{equation}\label{eq:projectiondegree}
\sum_v\tr(\Pi_v\mathcal F_v(s,A))
=\sum_vr_v\rk(\Pi_v)(\kappa_v-\kappa_E)
=-\sum_v\tr(\Pi_vb_vI).
\end{equation}
Conjugation by $e^{u/2}$ does not change this trace because it
preserves the summand. The moment-map trace on that summand is
zero: every remaining arrow begins and ends within the same
summand. Thus $\langle\Pi,\mathcal R(s,A)\rangle=0$.

The unitary $2\Pi-I$ fixes $A$ and reverses $C$. It commutes with
$\Pi$, so $\langle\Pi,\mathcal R(s,A+C)\rangle$ is even in $C$.
Its first derivative at $C=0$ vanishes. Taylor expansion in $C$
alone, and the Hessian bound along $A+tC$, now give
\[
|\langle\Pi,\mathcal R(s,A+C)\rangle|
\le Cq_s(A+C)q_s(C).
\]
Here $q_s(A+tC)\le q_s(A+C)$ for $0\le t\le1$ by orthogonality.
The finitely many rank choices of $\Pi$ give compact products of
Grassmannians; unitary invariance and bounded derivatives make
the constant independent of $\Pi$. This proves
\eqref{eq:projectionestimate}.
\end{proof}

Only the auxiliary representation $A$, with its crossing blocks removed,
was used to form the holomorphic direct summand in the proof.
The projection $\Pi$ need not define a holomorphic subbundle for the
original representation $Z$. This is why
\eqref{eq:projectionestimate} is available for every orthogonal
projection, including the spectral projections used next.
